\documentclass[a4paper,12pt]{article}

\usepackage[T2A]{fontenc}
\usepackage[utf8]{inputenc}
\usepackage[russian]{babel}

\usepackage{amsmath,amssymb,amsthm,mathtools}
\usepackage{helvet}
\renewcommand{\familydefault}{\sfdefault}

\usepackage{icomma}
\usepackage[a4paper,margin=2.2cm]{geometry}
\usepackage{microtype}

\usepackage{tikz}
\usetikzlibrary{calc}
\usepackage{tkz-euclide}

\usepackage{pgfplots}
\pgfplotsset{compat=1.18}

\usepackage{tabularx,booktabs,longtable}
\usepackage{enumitem}
\usepackage{hyperref}
\usepackage{parskip}
\usepackage{fancyhdr}
\usepackage{setspace}
\usepackage{xcolor}

% ------------------------------------------------
% Палитра
% ------------------------------------------------

\definecolor{accent}{HTML}{008080}
\definecolor{darkaccent}{HTML}{004D4D}
\definecolor{textgray}{HTML}{333333}
\definecolor{softgray}{HTML}{6B7474}
\definecolor{lightaccent}{HTML}{EAF4F4}

\hypersetup{
    colorlinks=true,
    linkcolor=darkaccent,
    urlcolor=darkaccent
}

\onehalfspacing
\setlength{\parindent}{0pt}
\setlength{\parskip}{0.55em}

\setlist[itemize]{
    leftmargin=1.5em,
    itemsep=0.3em,
    topsep=0.3em
}

\setlist[enumerate]{
    leftmargin=1.8em,
    itemsep=0.55em,
    topsep=0.4em
}

\pagestyle{fancy}
\fancyhf{}
\fancyhead[L]{\small\color{softgray} Теория вероятностей}
\fancyhead[R]{\small\color{softgray} Ксения Клыкова}
\fancyfoot[C]{\small\color{softgray}\thepage}
\renewcommand{\headrulewidth}{0.3pt}
\renewcommand{\headrule}{%
    \hbox to\headwidth{\color{lightaccent}\leaders\hrule height \headrulewidth\hfill}
}

\newcommand{\accentline}{%
    \par\noindent
    {\color{accent}\rule{\linewidth}{0.7pt}}
    \par
}

\newcommand{\key}[1]{\textbf{\color{darkaccent}#1}}

\begin{document}
\color{textgray}

% ================================================================
% Титульный лист
% ================================================================

\begin{titlepage}
\thispagestyle{empty}

\begin{center}

\vspace*{3.2cm}

{\Huge\bfseries Чек-лист}

\vspace{0.8cm}

{\LARGE\bfseries\color{darkaccent}
Теория вероятностей}

\vspace{0.35cm}

{\Large
Независимые испытания, формула Бернулли\\
и событие «хотя бы одно»}

\vspace{2.0cm}

{\large
Лёвин Артём Александрович\\
эксклюзивно для Ксении Клыковой}

\vspace{0.65cm}

{\normalsize ЕГЭ по математике}

\vspace{1.2cm}

{\large \today}

\end{center}

\vfill

\begin{center}
\begin{tikzpicture}[x=1cm,y=1cm]
    \draw[
        color=accent,
        line width=1.1pt
    ]
    (0,0)
    .. controls (3.0,0.65) and (7.3,-0.55) ..
    (10.5,0.05)
    .. controls (12.0,0.35) and (13.3,0.20) ..
    (14.2,0.05);
\end{tikzpicture}
\end{center}

\vspace{0.5cm}

\end{titlepage}

% ================================================================
% Основной чек-лист
% ================================================================

\section*{\color{darkaccent}1. Главная идея занятия}

В задачах с повторными испытаниями важно определить:

\begin{itemize}
    \item сколько всего проводится испытаний;
    \item какое событие считается нужным;
    \item какова вероятность этого события в одном испытании;
    \item требуется ли конкретное расположение успехов или только их количество;
    \item можно ли использовать противоположное событие для сокращения вычислений.
\end{itemize}

\accentline

\section*{\color{darkaccent}2. Союз «и»: вероятности перемножаем}

Если события независимы, то вероятность их совместного наступления равна произведению вероятностей:

\[
P(A\cap B)=P(A)\cdot P(B).
\]

Для нескольких независимых событий:

\[
P(A_1\cap A_2\cap\ldots\cap A_n)
=
P(A_1)P(A_2)\ldots P(A_n).
\]

\key{Смысл:} если в условии требуется, чтобы произошло первое событие \emph{и} второе \emph{и} третье, при независимости испытаний возникают множители.

\subsection*{Пример: пятёрка выпадает при каждом из шести бросков}

Вероятность выпадения пятёрки при одном броске:

\[
p=\frac16.
\]

Требуется:

\[
\text{пятёрка в 1-й раз}
\quad\text{и}\quad
\text{пятёрка во 2-й раз}
\quad\text{и}\quad \ldots
\quad\text{и}\quad
\text{пятёрка в 6-й раз}.
\]

Поэтому

\[
P=\left(\frac16\right)^6.
\]

\subsection*{Пример: пятёрка не выпадает ни разу}

Вероятность того, что при одном броске выпадет любое число, кроме пятёрки:

\[
q=1-\frac16=\frac56.
\]

Следовательно,

\[
P=\left(\frac56\right)^6.
\]

\accentline

\section*{\color{darkaccent}3. Конкретная позиция успеха}

Пусть пятёрка должна выпасть \key{только при первом броске}.

Тогда:

\[
\underbrace{\frac16}_{\text{выпала}}
\cdot
\underbrace{\frac56}_{\text{не выпала}}
\cdot
\underbrace{\frac56}_{\text{не выпала}}
\cdot
\underbrace{\frac56}_{\text{не выпала}}
\cdot
\underbrace{\frac56}_{\text{не выпала}}
\cdot
\underbrace{\frac56}_{\text{не выпала}}.
\]

То есть

\[
P=\frac16\left(\frac56\right)^5.
\]

Если пятёрка должна выпасть только при втором броске, получится

\[
P=
\frac56\cdot\frac16\cdot
\left(\frac56\right)^4.
\]

После перестановки множителей:

\[
P=\frac16\left(\frac56\right)^5.
\]

\key{Вывод:} вероятность одинакова для любой заранее указанной позиции единственного успеха.

\section*{\color{darkaccent}4. Союз «или»: вероятности вариантов складываем}

Если подходят несколько \key{несовместимых} вариантов, вероятности этих вариантов складываются.

Например, пятёрка может выпасть ровно один раз:

\[
\begin{aligned}
&\text{только в 1-м броске},\\
&\text{или только во 2-м},\\
&\text{или только в 3-м},\\
&\ldots\\
&\text{или только в 6-м}.
\end{aligned}
\]

Каждый из шести вариантов имеет вероятность

\[
\frac16\left(\frac56\right)^5.
\]

Поэтому

\[
P=
6\cdot
\frac16
\left(\frac56\right)^5
=
\left(\frac56\right)^5.
\]

\medskip

В общем случае для несовместимых событий

\[
P(A\cup B)=P(A)+P(B).
\]

Если события могут произойти одновременно, используется формула

\[
P(A\cup B)
=
P(A)+P(B)-P(A\cap B).
\]

\accentline

\section*{\color{darkaccent}5. Почему появляется формула Бернулли}

Если требуется ровно один успех, число возможных позиций невелико.

Если требуется ровно два успеха, приходится учитывать варианты:

\[
++----,\qquad +-+---,\qquad +--+--,\qquad \ldots
\]

При большом числе испытаний прямой перебор становится неудобным.

Для таких задач используется \key{формула Бернулли}.

\[
\boxed{
P_n(k)
=
\binom{n}{k}
p^k(1-p)^{n-k}
}
\]

где

\[
\binom{n}{k}
=
\frac{n!}{k!(n-k)!}.
\]

Следовательно,

\[
\boxed{
P_n(k)
=
\frac{n!}{k!(n-k)!}
p^k(1-p)^{n-k}
}
\]

\section*{\color{darkaccent}6. Что означает каждая буква}

\renewcommand{\arraystretch}{1.35}

\begin{table}[h!]
\centering
\caption{Обозначения в схеме Бернулли}
\begin{tabularx}{\linewidth}{@{}cX@{}}
\toprule
\textbf{Обозначение} & \textbf{Смысл} \\
\midrule
$n$ & общее количество независимых испытаний; \\
$k$ & нужное количество успешных испытаний; \\
$p$ & вероятность успеха в одном испытании; \\
$1-p$ & вероятность противоположного результата в одном испытании; \\
$\binom{n}{k}$ & количество способов расположить $k$ успехов среди $n$ испытаний. \\
\bottomrule
\end{tabularx}
\end{table}

\key{Ключевое условие схемы Бернулли:}

\begin{itemize}
    \item испытания независимы;
    \item вероятность успеха $p$ одинакова в каждом испытании;
    \item требуется определённое количество успехов из общего числа испытаний.
\end{itemize}

\accentline

\section*{\color{darkaccent}7. Что считается «успехом»}

Слово \emph{успех} в формуле Бернулли означает событие, которое требуется в конкретной задаче.

Оно может иметь любой смысл:

\begin{itemize}
    \item выпадение пятёрки;
    \item невыпадение пятёрки;
    \item попадание стрелка;
    \item промах стрелка;
    \item окончание кофе;
    \item сохранение кофе;
    \item выпадение решки.
\end{itemize}

\key{Всегда ориентируйтесь на вопрос задачи.}

Например, если

\[
P(\text{кофе останется})=0{,}2,
\]

то

\[
P(\text{кофе закончится})
=
1-0{,}2
=
0{,}8.
\]

Если требуется, чтобы кофе закончилось ровно в двух машинах из трёх, именно событие

\[
\text{«кофе закончилось»}
\]

следует считать успехом, поэтому

\[
n=3,\qquad k=2,\qquad p=0{,}8.
\]

\section*{\color{darkaccent}8. Пример применения формулы Бернулли}

Игральный кубик бросают шесть раз. Найдём вероятность того, что пятёрка выпадет ровно один раз.

Имеем:

\[
n=6,
\qquad
k=1,
\qquad
p=\frac16.
\]

По формуле Бернулли:

\[
P_6(1)
=
\frac{6!}{1!\,5!}
\left(\frac16\right)
\left(\frac56\right)^5.
\]

Так как

\[
\frac{6!}{5!}=6,
\]

то

\[
P_6(1)
=
6\cdot\frac16
\left(\frac56\right)^5
=
\left(\frac56\right)^5.
\]

\accentline

\section*{\color{darkaccent}9. Факториал}

Факториалом натурального числа $n$ называется произведение всех натуральных чисел от $1$ до $n$:

\[
n!=n(n-1)(n-2)\ldots 2\cdot1.
\]

Например,

\[
5!=5\cdot4\cdot3\cdot2\cdot1,
\]

\[
3!=3\cdot2\cdot1,
\]

\[
10!=10\cdot9\cdot8\cdots2\cdot1.
\]

Также

\[
0!=1.
\]

\subsection*{Полезное сокращение}

Необязательно каждый раз полностью раскрывать факториалы.

Например,

\[
10!
=
10\cdot9\cdot8!.
\]

Поэтому

\[
\frac{10!}{8!}
=
\frac{10\cdot9\cdot8!}{8!}
=
90.
\]

Аналогично,

\[
\frac{6!}{2!\,4!}
=
\frac{6\cdot5\cdot4!}{2\cdot1\cdot4!}
=
15.
\]

\key{Практический приём:} раскрывайте больший факториал только до меньшего факториала, который можно сократить.

\section*{\color{darkaccent}10. Алгоритм решения по формуле Бернулли}

\begin{enumerate}
    \item Определите, какое событие считается успехом.
    \item Найдите вероятность одного успеха $p$.
    \item Найдите противоположную вероятность:
    \[
    q=1-p.
    \]
    \item Определите общее число испытаний $n$.
    \item Определите нужное число успехов $k$.
    \item Подставьте данные:
    \[
    P_n(k)
    =
    \binom{n}{k}p^kq^{n-k}.
    \]
    \item Аккуратно сократите факториалы.
    \item Упростите степени и числовые множители.
\end{enumerate}

\accentline

\section*{\color{darkaccent}11. Событие «хотя бы один»}

Формулировка

\[
\text{«хотя бы один»}
\]

означает:

\[
1,\ 2,\ 3,\ \ldots
\]

успешных исходов.

Например, среди двух кофемашин выражение

\[
\text{«кофе закончится хотя бы в одной»}
\]

включает три варианта:

\[
\begin{array}{c}
\text{закончится только в первой},\\
\text{закончится только во второй},\\
\text{закончится в обеих}.
\end{array}
\]

Все эти случаи можно вычислять и складывать, однако обычно быстрее использовать противоположное событие.

\subsection*{Главная формула}

Событию

\[
\text{«хотя бы один успех»}
\]

противоположно событие

\[
\text{«ни одного успеха»}.
\]

Поэтому

\[
\boxed{
P(\text{хотя бы один успех})
=
1-P(\text{ни одного успеха})
}
\]

Если вероятность успеха в одном испытании равна $p$, то вероятность неуспеха равна $1-p$, поэтому

\[
\boxed{
P(\text{хотя бы один успех})
=
1-(1-p)^n
}
\]

\subsection*{Симметричная идея}

Если требуется

\[
\text{«хотя бы одно нет»},
\]

то противоположным событием является

\[
\text{«все да»}.
\]

Следовательно,

\[
\boxed{
P(\text{хотя бы одно нет})
=
1-P(\text{все да})
}
\]

А для события

\[
\text{«хотя бы одно да»}
\]

получаем

\[
\boxed{
P(\text{хотя бы одно да})
=
1-P(\text{все нет})
}
\]

\section*{\color{darkaccent}12. Полная группа событий}

Полной группой называется совокупность несовместимых событий, которая описывает все возможные исходы рассматриваемого опыта.

Сумма вероятностей событий полной группы равна единице:

\[
P(A_1)+P(A_2)+\ldots+P(A_m)=1.
\]

Именно это свойство лежит в основе перехода к противоположному событию:

\[
P(\overline A)=1-P(A).
\]

\accentline

\section*{\color{darkaccent}13. Типичные ошибки}

\begin{enumerate}
    \item \key{Подменить вероятность одного успеха долей нужных успехов.}

    Если монету бросают $10$ раз и требуется ровно $3$ решки, то

    \[
    p=\frac12,
    \]

    поскольку $p$ относится к одному броску.

    Число

    \[
    \frac3{10}
    \]

    вероятностью успеха в одном испытании не является.

    \item \key{Неверно выбрать успех.}

    Если дано

    \[
    P(\text{кофе останется})=0{,}2,
    \]

    а требуется окончание кофе, то

    \[
    p=0{,}8.
    \]

    \item \key{Забыть коэффициент $\binom nk$.}

    Выражение

    \[
    p^k(1-p)^{n-k}
    \]

    описывает только одно конкретное расположение успехов и неуспехов.

    \item \key{Использовать формулу Бернулли при разных вероятностях успеха.}

    В стандартной схеме Бернулли вероятность $p$ должна оставаться одинаковой.

    \item \key{Механически применять правило «и --- умножаем».}

    Произведение вероятностей напрямую используется для независимых событий.

    \item \key{Механически применять правило «или --- складываем».}

    Простая сумма подходит для несовместимых вариантов.

    \item \key{Считать все случаи для «хотя бы одного», когда проще взять дополнение.}

    Обычно выгоднее:

    \[
    P(\text{хотя бы один})
    =
    1-P(\text{ни одного}).
    \]
\end{enumerate}

\section*{\color{darkaccent}14. Мини-чек перед решением}

Перед вычислениями проверьте:

\begin{itemize}
    \item[$\square$] Что является одним испытанием?
    \item[$\square$] Какое событие требуется считать успехом?
    \item[$\square$] Чему равно $p$?
    \item[$\square$] Чему равно $1-p$?
    \item[$\square$] Сколько всего испытаний $n$?
    \item[$\square$] Сколько успехов требуется $k$?
    \item[$\square$] Испытания независимы?
    \item[$\square$] Вероятность $p$ одинакова во всех испытаниях?
    \item[$\square$] Требуется конкретное расположение успехов или только их число?
    \item[$\square$] Есть ли слова «хотя бы один», позволяющие использовать противоположное событие?
\end{itemize}

% ================================================================
% Тренировочный блок
% ================================================================

\newpage

\section*{\color{darkaccent}Тренировочный блок}

\textbf{Формат: Путь А}

Во всех задачах испытания считать независимыми.

\subsection*{\color{darkaccent}Средний уровень}

\begin{enumerate}[label=\textbf{\arabic*.}]

    \item Игральный кубик бросают шесть раз. Найдите вероятность того, что пятёрка выпадет при каждом броске.

    \item Игральный кубик бросают шесть раз. Найдите вероятность того, что пятёрка не выпадет ни разу.

    \item Игральный кубик бросают шесть раз. Найдите вероятность того, что пятёрка выпадет ровно один раз.

\end{enumerate}

\subsection*{\color{darkaccent}Выше среднего}

\begin{enumerate}[label=\textbf{\arabic*.},resume]

    \item Игральный кубик бросают шесть раз. Найдите вероятность того, что пятёрка выпадет ровно два раза.

    \item Монету бросают десять раз. Найдите вероятность того, что решка выпадет ровно три раза.

    \item Стрелок производит шесть независимых выстрелов. Вероятность попадания при каждом выстреле равна $0{,}8$. Найдите вероятность ровно четырёх попаданий.

    \item Имеются три кофемашины. Вероятность того, что кофе останется в каждой отдельной машине, равна $0{,}2$. Найдите вероятность того, что кофе закончится ровно в двух машинах.

    \item Игральный кубик бросают шесть раз. Найдите вероятность того, что пятёрка выпадет хотя бы один раз.

    \item Монету бросают двенадцать раз. Найдите вероятность того, что орёл выпадет не менее десяти раз.

\end{enumerate}

\subsection*{\color{darkaccent}Сложная задача}

\begin{enumerate}[label=\textbf{\arabic*.},resume]

    \item Вероятность того, что случайно выбранное изделие имеет дефект, равна $0{,}1$. Проверяют $12$ изделий. Найдите вероятность того, что дефектными окажутся как минимум два, но не более четырёх изделий.

\end{enumerate}

% ================================================================
% Ответы
% ================================================================

\newpage

\section*{\color{darkaccent}Ответы к тренировочному блоку}

\renewcommand{\arraystretch}{1.5}

\begin{table}[h!]
\centering
\caption{Ответы}
\begin{tabularx}{\linewidth}{@{}cX@{}}
\toprule
\textbf{№} & \textbf{Ответ} \\
\midrule

1 &
\[
\left(\frac16\right)^6
=
\frac1{46656}.
\]
\\

2 &
\[
\left(\frac56\right)^6
=
\frac{15625}{46656}.
\]
\\

3 &
\[
\binom61\frac16\left(\frac56\right)^5
=
\left(\frac56\right)^5
=
\frac{3125}{7776}.
\]
\\

4 &
\[
\binom62
\left(\frac16\right)^2
\left(\frac56\right)^4
=
\frac{3125}{15552}.
\]
\\

5 &
\[
\binom{10}{3}
\left(\frac12\right)^{10}
=
\frac{15}{128}.
\]
\\

6 &
\[
\binom64
(0{,}8)^4(0{,}2)^2
=
0{,}24576.
\]
\\

7 &
\[
\binom32
(0{,}8)^2(0{,}2)
=
0{,}384.
\]
\\

8 &
\[
1-\left(\frac56\right)^6
=
\frac{31031}{46656}.
\]
\\

9 &
\[
\binom{12}{10}
\left(\frac12\right)^{12}
+
\binom{12}{11}
\left(\frac12\right)^{12}
+
\binom{12}{12}
\left(\frac12\right)^{12}
=
\frac{79}{4096}.
\]
\\

10 &
\[
\binom{12}{2}(0{,}1)^2(0{,}9)^{10}
+
\binom{12}{3}(0{,}1)^3(0{,}9)^9
+
\binom{12}{4}(0{,}1)^4(0{,}9)^8.
\]
\\

\bottomrule
\end{tabularx}
\end{table}

\vfill

\begin{center}
{\small\color{softgray}
Главный ориентир: сначала определить нужное событие и вероятность одного испытания,\\
затем выбрать между произведением, суммой вариантов, формулой Бернулли\\
и переходом к противоположному событию.}
\end{center}

\end{document}